\chapter{Organising and Protecting Your Data}
\label{ch:organisation}

This chapter is about filing and backups. It is in a book about AI because
agentic workflows amplify both the value of good organisation and the cost of
bad organisation, and because an agent with write access to a badly-backed-up
directory is a new category of risk.

\section{Why this matters more now}

An agent navigates your project by reading it. A directory tree with meaningful
names, a \code{README.md} that explains the structure, and consistent file
naming lets an agent locate the right file in one or two operations. A tree
named \code{final/}, \code{final2/}, \code{new\_final/}, \code{Untitled1.csv}
forces it to read everything, burn context on irrelevant material, and guess.

The second reason is blunter. Agents delete files. Sometimes because you asked
them to, sometimes through a misread instruction, occasionally through a bug.
\textcite{shapira2026agentsofchaos} document execution of destructive
system-level actions among the failure modes observed when autonomous agents
were given real file system and shell access in a live environment. Your backup
strategy is the thing standing between that and a lost field season.

\section{PARA}

PARA is a general-purpose organisation method that adapts well to research. Four
top-level categories, by \emph{actionability} rather than topic:

\begin{description}
  \item[Projects] Active work with a defined outcome and an end date. A
    manuscript in preparation, a specific analysis, a deliverable.
  \item[Areas] Ongoing responsibilities with no end date. A long-term monitoring
    programme, a database you maintain, teaching.
  \item[Resources] Reference material of future interest. Literature,
    protocols, reusable code, notes.
  \item[Archive] Anything from the first three that is finished or dormant.
\end{description}

The mechanism that makes it work is the archive. When a project closes it moves
to \code{Archive/} intact. The active tree stays small, which keeps it
navigable --- for you and for an agent.

\begin{lstlisting}[style=shell,caption={A research tree, PARA-shaped}]
~/work/
  10_projects/
    2026_posidonia_edna/
      README.md          # what, why, status, how to run
      AGENTS.md          # conventions for agents working here
      data/
        raw/             # never edited, read-only, checksummed
        processed/       # derived; regenerable from raw + scripts
      scripts/
      results/
      manuscript/
  20_areas/
    monitoring_network/
    teaching/
  30_resources/
    literature/
    protocols/
  40_archive/
    2024_thermal_tolerance/
\end{lstlisting}

The numeric prefixes keep sort order stable and make paths easy to type and to
reference in conversation.

\section{Naming}

Three rules cover most of it.

\textbf{Sortable dates.} \code{YYYY-MM-DD} or \code{YYYYMMDD}. These sort
lexicographically into chronological order, which no other format does.

\textbf{No spaces or special characters.} Underscores between fields, hyphens
within them. Spaces break shell pipelines, require quoting everywhere, and
produce a steady trickle of bugs in scripts written by humans and agents alike.

\textbf{Machine-readable structure.} Encode the metadata in the name, in
consistent field order:

\begin{lstlisting}[style=shell]
20260415_crete-north_station03_16S_rep1.fastq.gz
20260415_crete-north_station03_16S_rep2.fastq.gz
^date     ^site        ^station  ^marker ^replicate
\end{lstlisting}

This is greppable, globbable, and parseable into a data frame with one split.

A name a script can select and split is worth more than a descriptive sentence.
Put the description in the README.

\section{Raw data is read-only}

The single most valuable convention in this chapter: raw data is never modified.
Set it read-only, checksum it, and derive everything else.

\begin{lstlisting}[style=shell]
chmod -R a-w data/raw/
find data/raw -type f -exec sha256sum {} + > data/raw/CHECKSUMS.sha256
sha256sum -c data/raw/CHECKSUMS.sha256      # verify later
\end{lstlisting}

Everything in \code{processed/} should be reconstructible by running the scripts
against \code{raw/}. If it is not, that is a gap in the pipeline, and it is
better to find out now than during review.

This convention also protects you from agents. An agent instructed to clean a
dataset will happily clean it in place. Read-only permissions turn that into a
visible error rather than a silent loss.

\section{Backup: the 3-2-1 rule}

Three copies of the data, on two different media, with one off-site.

\begin{itemize}
  \item \textbf{Copy 1} --- the working copy on your machine.
  \item \textbf{Copy 2} --- a local backup on different physical media. An
    external drive, a NAS. Different media matters: a second partition on the
    same disk dies with the disk.
  \item \textbf{Copy 3} --- off-site. Institutional storage, cloud sync, a drive
    at another building. This is the copy that survives fire, theft and flood.
\end{itemize}

\subsection{RAID is not a backup}

RAID protects against \emph{drive failure}. It does not protect against
accidental deletion, ransomware, filesystem corruption, controller failure, a
power event, theft, or an agent running \code{rm -rf} against the wrong path.
Every one of those propagates instantly to every disk in the array.

A NAS with RAID is a good copy~2. It is not copies 2 and 3, and treating it as
both is the most common backup mistake in research computing.

\subsection{Verify restores}

An unverified backup is a hypothesis. Schedule a restore test --- pick a random
file, restore it from each copy, diff it against the original. Quarterly is
enough. The failure you are testing for is the one where backups have been
silently failing for eight months.

\section{Version control}

Git for code, always. For a project that involves agents, it stops being
optional: \code{git diff} is how you review what an agent actually changed,
\code{git checkout} is how you undo it, and commit granularity is how you
bisect when something breaks three steps later.

\begin{lstlisting}[style=shell,caption={Minimum viable discipline}]
git init
git add README.md AGENTS.md scripts/ .gitignore
git commit -m "Initial structure"

# commit before letting an agent loose
git commit -am "checkpoint before refactor"
# ... agent works ...
git diff                    # review every line
git checkout -- .           # or undo the lot
\end{lstlisting}

Large data files do not belong in Git. Use Git LFS, DVC, or --- most often the
right answer --- keep data out of the repository entirely and record its location
and checksum in the README. Zenodo gives you a DOI and long-term storage for the
archival copy.

\section{Plain text lasts}

A closing principle that runs through the rest of this book. The formats that
survive are the simple ones: plain text, Markdown, CSV, TSV, and well-specified
binary formats with open readers such as Parquet and NetCDF.

Plain text is greppable, diffable, version-controllable, readable by every tool,
and legible in fifty years without a specific application. It is also the format
agents work with best, which is not a coincidence --- they were trained on it.

Proprietary formats, database-internal state and application-specific exports are
where data goes to become inaccessible. When a system offers you a choice, take
the text.
